%!TeX program = xelatex
\documentclass{SYSUReport}

% 根据个人情况修改
\headl{实验二：中断处理}
\headr{操作系统原理实验}
\lessonTitle{操作系统原理实验}
\reportTitle{实验二：中断处理}
\inst{计算机学院}
\major{保密管理}
\stuname{吕梓文}
\stuid{24353028}
\date{\today}

\AddEnumerateCounter{\chinese}{\chinese}{}
\begin{document}

\cover
\updateheader{实验二：中断处理}
\thispagestyle{empty} % 首页不显示页码
\clearpage

% \pagenumbering{Roman}
% \setcounter{page}{1}
% \tableofcontents
\clearpage

\pagenumbering{arabic}
\setcounter{page}{1}

\section{实验完成情况}

\begin{itemize}
    \item \textbf{必做实验内容}：全部完成
    \begin{itemize}
        \item GDT 与 TSS 实现（为 Double Fault 和 Page Fault 准备独立栈）
        \item 注册中断处理程序（IDT 初始化）
        \item 初始化 APIC（XAPIC 和 IOAPIC）
        \item 时钟中断实现
        \item 串口输入中断实现
        \item 用户交互（输入缓冲区和命令行）
    \end{itemize}
    \item \textbf{加分项}：完成 4 项中的 3 项（除第 2 项外）
    \begin{itemize}
        \item 为全部可能的 CPU 异常设置处理程序
        \item 使用 bitflags 和 bit\_field 定义 APIC 寄存器
        \item 串口输入驱动支持中文和 emoji 输入
    \end{itemize}
\end{itemize}

\section{GDT 与 TSS}

\begin{enumerate}[label={\chinese*、},labelsep=0pt]

    \item \textbf{TSS 中断栈表实现}

    在 x86\_64 架构中，TSS（Task State Segment）用于存储任务状态信息，其中 IST（Interrupt Stack Table）字段可以为特定的中断处理程序指定独立的栈空间。这是为了防止某些严重异常（如 Double Fault）发生时，由于当前栈已损坏而导致系统崩溃。

    在 \texttt{src/memory/gdt.rs} 中，首先定义了两个常量来标识 IST 中的索引位置：

    \begin{minted}{rust}
pub const DOUBLE_FAULT_IST_INDEX: u16 = 0;
pub const PAGE_FAULT_IST_INDEX: u16 = 1;

pub const IST_SIZES: [usize; 3] = [0x1000, 0x1000, 0x1000];
    \end{minted}

    \texttt{DOUBLE\_FAULT\_IST\_INDEX} 设为 0，表示 Double Fault 处理程序将使用 IST 的第一个条目；\texttt{PAGE\_FAULT\_IST\_INDEX} 设为 1，表示 Page Fault 处理程序使用第二个条目。每个 IST 栈的大小为 4KB（\texttt{0x1000} 字节），这个大小足以容纳异常处理过程中的函数调用栈。

    接下来使用 \texttt{lazy\_static} 宏定义全局静态的 TSS 结构体：

    \begin{minted}{rust}
lazy_static! {
    static ref TSS: TaskStateSegment = {
        let mut tss = TaskStateSegment::new();

        // ... privilege_stack_table 初始化 ...

        // Allocate IST stack for Double Fault
        tss.interrupt_stack_table[DOUBLE_FAULT_IST_INDEX as usize] = {
            const STACK_SIZE: usize = IST_SIZES[0];
            static mut STACK: [u8; STACK_SIZE] = [0; STACK_SIZE];
            let stack_start = VirtAddr::from_ptr(addr_of_mut!(STACK));
            let stack_end = stack_start + STACK_SIZE as u64;
            info!(
                "DOUBLE FAULT IST  : 0x{:016x}-0x{:016x}",
                stack_start.as_u64(),
                stack_end.as_u64()
            );
            stack_end
        };
        
        // Allocate IST stack for Page Fault
        tss.interrupt_stack_table[PAGE_FAULT_IST_INDEX as usize] = {
            const STACK_SIZE: usize = IST_SIZES[1];
            static mut STACK: [u8; STACK_SIZE] = [0; STACK_SIZE];
            let stack_start = VirtAddr::from_ptr(addr_of_mut!(STACK));
            let stack_end = stack_start + STACK_SIZE as u64;
            info!(
                "PAGE FAULT IST  : 0x{:016x}-0x{:016x}",
                stack_start.as_u64(),
                stack_end.as_u64()
            );
            stack_end
        };

        tss
    };
}
    \end{minted}

    这里需要注意几个关键点。首先，IST 栈使用静态缓冲区分配在 BSS 段，这是因为在内核启动早期，堆分配器尚未初始化，无法使用动态内存分配。其次，\texttt{tss.interrupt\_stack\_table} 存储的是栈顶地址而非栈底地址，这是因为 x86\_64 架构的栈是向下增长的，CPU 在切换栈时会直接使用这个值作为新的 RSP。

    使用 \texttt{addr\_of\_mut!} 宏获取静态变量的地址是 Rust 中处理可变静态变量的安全方式，它避免了直接引用可变静态变量可能导致的未定义行为。

    \item \textbf{GDT 初始化}

    GDT（Global Descriptor Table）是 x86 架构中用于定义内存段属性的数据结构。在现代 64 位操作系统中，虽然段机制已被弱化，但 GDT 仍然需要存在，因为它还包含了 TSS 描述符。

    \begin{minted}{rust}
lazy_static! {
    static ref GDT: (GlobalDescriptorTable, KernelSelectors) = {
        let mut gdt = GlobalDescriptorTable::new();
        let code_selector = gdt.append(Descriptor::kernel_code_segment());
        let data_selector = gdt.append(Descriptor::kernel_data_segment());
        let tss_selector = gdt.append(Descriptor::tss_segment(&TSS));
        (
            gdt,
            KernelSelectors {
                code_selector,
                data_selector,
                tss_selector,
            },
        )
    };
}
    \end{minted}

    GDT 中添加了三个描述符：内核代码段、内核数据段和 TSS 段。TSS 描述符是必需的，因为 CPU 需要通过 \texttt{LTR} 指令加载 TSS 选择子来激活 TSS。

    GDT 的加载和段寄存器的设置在 \texttt{init} 函数中完成：

    \begin{minted}{rust}
pub fn init() {
    use x86_64::{
        PrivilegeLevel,
        instructions::{
            segmentation::{CS, DS, ES, FS, GS, SS},
            tables::load_tss,
        },
    };

    GDT.0.load();
    unsafe {
        CS::set_reg(GDT.1.code_selector);
        DS::set_reg(GDT.1.data_selector);
        SS::set_reg(SegmentSelector::new(0, PrivilegeLevel::Ring0));
        ES::set_reg(SegmentSelector::new(0, PrivilegeLevel::Ring0));
        FS::set_reg(SegmentSelector::new(0, PrivilegeLevel::Ring0));
        GS::set_reg(SegmentSelector::new(0, PrivilegeLevel::Ring0));
        load_tss(GDT.1.tss_selector);
    }
    // ...
}
    \end{minted}

    这里需要按顺序执行：首先加载 GDT，然后更新各个段寄存器，最后加载 TSS。段寄存器的更新必须使用 \texttt{unsafe} 块，因为这会改变 CPU 的运行状态。需要注意的是，在 64 位模式下，CS、DS 等段寄存器的基址和限长实际上被忽略，但它们的值仍然需要正确设置以维持 CPU 的正常运行状态。

    \begin{figure}[h]
        \centering
        \includegraphics[width=0.9\textwidth]{local_apic_register.png}
        \caption{内核启动输出 - GDT 和 IST 初始化}
    \end{figure}

    \item \textbf{思考题：TSS 栈空间不足的情况}

    \textbf{【思考题 6】} 如果在 TSS 中为中断分配的栈空间不足，会发生什么情况？

    当 IST 栈空间不足时，中断处理程序在执行过程中可能会发生栈溢出。由于 IST 栈是固定大小的静态缓冲区，如果中断处理程序需要更多的栈空间（例如深层函数调用、大量局部变量），就会覆盖相邻内存区域，导致数据损坏或系统崩溃。

    在 CPU 异常发生过程中：
    \begin{enumerate}
        \item CPU 检测到异常，查找 IDT 中对应的处理程序
        \item 如果处理程序设置了 \texttt{set\_stack\_index}，CPU 会切换到 IST 指定的栈
        \item 如果 IST 栈空间不足，在压入寄存器状态或执行处理程序时会发生栈溢出
        \item 栈溢出可能触发 Page Fault 或 General Protection Fault
        \item 如果此时正在处理 Double Fault，新的异常会导致 Triple Fault，系统重置
    \end{enumerate}

    Triple Fault 通常发生在以下情况：
    \begin{itemize}
        \item Double Fault 处理程序本身发生异常
        \item IDT 中 Double Fault 处理程序未正确设置
        \item IST 栈溢出导致 Double Fault 处理程序无法正常执行
    \end{itemize}

    在实际调试中，如果发现系统频繁重启且没有任何错误输出，很可能是发生了 Triple Fault。此时可以通过 QEMU 的 \texttt{-d int,cpu\_reset} 选项查看中断和重置的详细信息。

\end{enumerate}

\section{注册中断处理程序}

\begin{enumerate}[label={\chinese*、},labelsep=0pt, start=2]

    \item \textbf{IDT 初始化}

    IDT（Interrupt Descriptor Table）是 x86 架构中用于定义中断处理程序入口的数据结构。每个 IDT 条目包含处理程序的地址、段选择子和属性标志。

    在 \texttt{src/interrupt/mod.rs} 中，使用 \texttt{lazy\_static} 定义全局静态的 IDT，并将注册工作委托给各个子模块：

    \begin{minted}{rust}
mod apic;
mod consts;
pub mod clock;
mod serial;
mod exceptions;

use apic::*;
use x86_64::structures::idt::InterruptDescriptorTable;

use crate::{interrupt::consts::Irq, memory::physical_to_virtual};

lazy_static! {
    static ref IDT: InterruptDescriptorTable = {
        let mut idt = InterruptDescriptorTable::new();
        unsafe {
            exceptions::register_idt(&mut idt);
            clock::register_idt(&mut idt);
            serial::register_idt(&mut idt);
        }
        idt
    };
}
    \end{minted}

    这种模块化设计使得代码结构清晰，每个模块只负责注册自己关心的中断类型。\texttt{exceptions} 模块处理 CPU 异常，\texttt{clock} 模块处理时钟中断，\texttt{serial} 模块处理串口中断。

    \item \textbf{CPU 异常处理程序}

    在 \texttt{src/interrupt/exceptions.rs} 中，为各种 CPU 异常注册处理程序。x86\_64 架构定义了多种异常类型，分为 Fault（错误）、Trap（陷阱）和 Abort（终止）三类。

    \begin{minted}{rust}
pub unsafe fn register_idt(idt: &mut InterruptDescriptorTable) {
    // Faults
    idt.divide_error.set_handler_fn(divide_error_handler);
    idt.bound_range_exceeded.set_handler_fn(bound_range_exceeded_handler);
    idt.invalid_opcode.set_handler_fn(invalid_opcode_handler);
    idt.device_not_available.set_handler_fn(device_not_available_handler);
    idt.invalid_tss.set_handler_fn(invalid_tss_handler);
    idt.segment_not_present.set_handler_fn(segment_not_present_handler);
    idt.stack_segment_fault.set_handler_fn(stack_segment_fault_handler);
    idt.general_protection_fault.set_handler_fn(general_protection_fault_handler);
    unsafe {
        idt.page_fault
        .set_handler_fn(page_fault_handler)
        .set_stack_index(gdt::PAGE_FAULT_IST_INDEX);
    }
    idt.x87_floating_point.set_handler_fn(x87_floating_point_handler);
    idt.alignment_check.set_handler_fn(alignment_check_handler);
    idt.simd_floating_point.set_handler_fn(simd_floating_point_handler);

    // Traps
    idt.debug.set_handler_fn(debug_handler);
    idt.breakpoint.set_handler_fn(breakpoint_handler);
    idt.overflow.set_handler_fn(overflow_handler);

    // Aborts
    unsafe {
        idt.double_fault
        .set_handler_fn(double_fault_handler)
        .set_stack_index(gdt::DOUBLE_FAULT_IST_INDEX);
    }
    idt.machine_check.set_handler_fn(machine_check_handler);
}
    \end{minted}

    这里需要注意，Page Fault 和 Double Fault 的处理程序通过 \texttt{set\_stack\_index} 指定了独立的 IST 栈。这是因为这两种异常可能在系统状态已经受损的情况下发生，使用独立栈可以确保处理程序有足够的栈空间执行。

    中断处理函数使用 \texttt{extern "x86-interrupt"} 调用约定声明。这个调用约定是 Rust 专门为 x86 中断处理提供的，它会自动处理中断栈帧的保存和恢复，以及在中断返回时执行 \texttt{iretq} 指令。

    \begin{minted}{rust}
pub extern "x86-interrupt" fn divide_error_handler(stack_frame: InterruptStackFrame) {
    panic!("EXCEPTION: DIVIDE ERROR\n\n{:#?}", stack_frame);
}

pub extern "x86-interrupt" fn page_fault_handler(
    stack_frame: InterruptStackFrame,
    err_code: PageFaultErrorCode,
) {
    panic!(
        "EXCEPTION: PAGE FAULT, ERROR_CODE: {:?}\n\nTrying to access: {:#x}\n{:#?}",
        err_code,
        Cr2::read().unwrap_or(VirtAddr::new_truncate(0xdeadbeef)),
        stack_frame
    );
}
    \end{minted}

    Page Fault 处理程序会读取 Cr2 寄存器，该寄存器保存了导致 Page Fault 的虚拟地址。这对于调试非常有用，可以快速定位问题所在。

    \item \textbf{加分项：全部 CPU 异常处理程序}

    \textbf{【加分项 1】} 为全部可能的 CPU 异常设置对应的处理程序：

    如上述代码所示，已为以下异常注册了处理程序：
    \begin{itemize}
        \item \textbf{Faults}：Divide Error（除零错误）、Bound Range Exceeded（边界超限）、Invalid Opcode（无效操作码）、Device Not Available（设备不可用）、Invalid TSS（无效 TSS）、Segment Not Present（段不存在）、Stack Segment Fault（栈段错误）、General Protection Fault（一般保护错误）、Page Fault（页错误）、x87 Floating Point（x87 浮点异常）、Alignment Check（对齐检查）、SIMD Floating Point（SIMD 浮点异常）
        \item \textbf{Traps}：Debug（调试）、Breakpoint（断点）、Overflow（溢出）
        \item \textbf{Aborts}：Double Fault（双重错误）、Machine Check（机器检查）
    \end{itemize}

    每个处理程序都会输出异常类型和栈帧信息，方便调试。对于带有错误码的异常（如 General Protection Fault、Page Fault），处理程序还会输出错误码的具体值。

    \item \textbf{思考题：未使用 set\_stack\_index 时的栈位置}

    \textbf{【思考题 7】} 在未使用 \texttt{set\_stack\_index} 函数时，中断处理程序的栈可能在哪里？

    当未使用 \texttt{set\_stack\_index} 时，中断处理程序会使用当前 CPU 的栈指针（RSP）。具体来说：
    \begin{itemize}
        \item 如果中断发生在内核态（Ring 0），则使用当前内核栈
        \item 如果中断发生在用户态（Ring 3），CPU 会自动切换到 TSS 中的 RSP0（特权级栈）
    \end{itemize}

    在本实验中，由于所有代码都在内核态运行，未设置 \texttt{set\_stack\_index} 的中断处理程序会使用触发中断时正在使用的栈。这可能导致问题，因为如果中断发生在某个函数的深层调用中，栈空间可能不足。

    通过 GDB 调试器可以验证这一点：在中断处理程序中设置断点，查看 RSP 的值，可以发现它与触发中断时的栈地址相同。

    \item \textbf{思考题：Page Fault 属于可恢复异常的原因}

    \textbf{【思考题 5 - Page Fault 部分】} Page Fault 属于可恢复的异常，因为：

    Page Fault 的发生通常是因为访问的页面不在物理内存中（需要换入），而不是因为硬件错误或非法访问。Cr2 寄存器保存了导致 Page Fault 的虚拟地址，操作系统可以：
    \begin{enumerate}
        \item 检查该地址是否在进程的合法地址空间内
        \item 如果是合法的按需分页访问，从磁盘加载页面到物理内存
        \item 更新页表，重新执行导致 Page Fault 的指令
    \end{enumerate}

    从代码中可以看到，Page Fault 处理程序会读取 Cr2 寄存器：

    \begin{minted}{rust}
pub extern "x86-interrupt" fn page_fault_handler(
    stack_frame: InterruptStackFrame,
    err_code: PageFaultErrorCode,
) {
    panic!(
        "EXCEPTION: PAGE FAULT, ERROR_CODE: {:?}\n\nTrying to access: {:#x}\n{:#?}",
        err_code,
        Cr2::read().unwrap_or(VirtAddr::new_truncate(0xdeadbeef)),
        stack_frame
    );
}
    \end{minted}

    Cr2 寄存器的值指示了触发 Page Fault 的具体虚拟地址，这对于调试和实现虚拟内存管理非常重要。Page Fault 的错误码（\texttt{PageFaultErrorCode}）还包含了更多信息，如访问类型（读/写）、权限级别（用户/内核）等。

\end{enumerate}

\section{初始化 APIC}

\begin{enumerate}[label={\chinese*、},labelsep=0pt, start=4]

    \item \textbf{XAPIC 实现}

    APIC（Advanced Programmable Interrupt Controller）是 x86 架构中的高级可编程中断控制器，负责管理和分发中断。XAPIC 是 APIC 的一种实现方式，通过 MMIO（Memory-Mapped I/O）进行访问。

    在 \texttt{src/interrupt/apic/xapic.rs} 中，首先定义了 XAPIC 的结构体和基本的读写操作：

    \begin{minted}{rust}
pub struct XApic {
    addr: u64,
}

impl XApic {
    pub unsafe fn new(addr: u64) -> Self {
        XApic { addr }
    }

    unsafe fn read(&self, reg: u32) -> u32 {
        unsafe { read_volatile((self.addr + reg as u64) as *const u32) }
    }

    unsafe fn write(&mut self, reg: u32, value: u32) {
        unsafe {
            write_volatile((self.addr + reg as u64) as *mut u32, value);
            self.read(0x20);
        }
    }
}
    \end{minted}

    \texttt{read} 和 \texttt{write} 函数使用 \texttt{read\_volatile} 和 \texttt{write\_volatile} 进行 MMIO 操作。这些函数确保编译器不会优化掉对硬件寄存器的访问。

    \texttt{write} 函数中的 \texttt{self.read(0x20)} 是一个重要的技巧：读取 APIC ID 寄存器（偏移量 0x20）用于确保写操作完成。这是因为 MMIO 写操作可能是异步的，CPU 可能在实际写入完成前继续执行后续指令。通过读取一个无关寄存器，可以强制 CPU 等待写操作完成，这称为“序列化”操作。

    \item \textbf{APIC 初始化}

    \texttt{LocalApic} trait 的 \texttt{cpu\_init} 方法实现了完整的 APIC 初始化流程：

    \begin{minted}{rust}
impl LocalApic for XApic {
    fn support() -> bool {
        CpuId::new().get_feature_info().map(
            |f| f.has_apic()
        ).unwrap_or(false)
    }

    fn cpu_init(&mut self) {
        unsafe {
            // Enable local APIC; set spurious interrupt vector.
            let mut sivr = self.read(0xF0);
            sivr |= LAPICFlags::ENAPIC.bits();
            sivr &= !(0xFF);
            sivr |= Interrupts::IrqBase as u32 + Irq::Spurious as u32;
            self.write(0xF0, sivr);

            // Set LVT timer register
            let mut lvt_timer = self.read(0x320);
            lvt_timer &= !(0xFF);
            lvt_timer |= Interrupts::IrqBase as u32 + Irq::Timer as u32;
            lvt_timer &= LVTTimerFlags::CLEAR_MASK.bits();
            lvt_timer |= LVTTimerFlags::TIMER_PERIODIC_MODE.bits();
            self.write(0x320, lvt_timer);

            // Set TDCR and TICR
            self.write(0x3E0, TDCRFlags::TIMEDIVIDEBY1.bits());
            self.write(0x380, 0x20000);

            // Disable LINT0, LINT1, PCINT
            self.write(0x350, !LVTTimerFlags::CLEAR_MASK.bits());
            self.write(0x360, !LVTTimerFlags::CLEAR_MASK.bits());
            self.write(0x340, !LVTTimerFlags::CLEAR_MASK.bits());
            
            // Map error interrupt
            let mut lvt_error = self.read(0x370);
            lvt_error &= !(0xFF);
            lvt_error |= Interrupts::IrqBase as u32 + Irq::Error as u32;
            lvt_error &= LVTTimerFlags::CLEAR_MASK.bits();
            self.write(0x370, lvt_error);

            // Clear error status register
            self.write(0x280, 0);
            self.write(0x280, 0);

            // Ack outstanding interrupts
            self.write(0x0B, 0);

            // Send Init Level De-Assert
            self.write(0x310, 0);
            const BCAST: u32 = 1 << 19;
            const INIT: u32 = 5 << 8;
            const TMLV: u32 = 1 << 15;
            self.write(0x300, BCAST | INIT | TMLV);
            const DS: u32 = 1 << 12;
            while self.read(0x300) & DS != 0 {}

            // Enable interrupts
            self.write(0x08, 0);
        }
    }
    // ...
}
    \end{minted}

    APIC 初始化需要按照特定顺序执行多个步骤：

    首先，通过 CPUID 指令检查系统是否支持 APIC。这是必要的，因为某些旧系统或虚拟机可能不支持 APIC。

    然后，配置 SPIV（Spurious Interrupt Vector）寄存器（偏移量 0xF0）来启用 APIC。这个寄存器的第 8 位（EN 位）控制 APIC 的启用状态，低 8 位设置伪中断向量号。需要注意的是，SPIV 寄存器的低 4 位是只读的，始终为 1，因此在修改时需要保留这些位。

    接下来，配置 LVT（Local Vector Table）Timer 寄存器（偏移量 0x320）来设置定时器中断。需要设置中断向量号、清除屏蔽位（Mask），并设置周期模式（Periodic Mode）。周期模式下，定时器会在每次计数到 0 时自动重新加载初始值并触发中断。

    TDCR（Timer Divide Configuration Register，偏移量 0x3E0）设置分频系数，TICR（Timer Initial Count Register，偏移量 0x380）设置初始计数值。这两个寄存器共同决定了定时器中断的频率。

    禁用 LINT0、LINT1 和 PCINT 是为了防止这些中断线产生意外的中断。这些中断线通常用于连接外部设备，但在本实验中不需要使用。

    清除错误状态寄存器需要连续写入两次 0，这是 APIC 的硬件设计要求。

    发送 INIT De-assert IPI（Inter-Processor Interrupt）是为了同步 APIC 的仲裁 ID。虽然本实验是单核系统，但这个步骤仍然是初始化流程的一部分。

    \item \textbf{加分项：使用 bitflags 定义 APIC 寄存器}

    \textbf{【加分项 3】} 使用 \texttt{bitflags} 宏定义 APIC 寄存器标志位，提高代码可读性：

    \begin{minted}{rust}
bitflags! {
    #[derive(Debug, Clone, Copy, PartialEq, Eq, PartialOrd, Ord, Hash)]
    pub struct LAPICFlags: u32 {
        const ENAPIC = 1 << 8;
    }

    #[derive(Debug, Clone, Copy, PartialEq, Eq, PartialOrd, Ord, Hash)]
    pub struct LVTTimerFlags: u32 {
        const CLEAR_MASK = !(1 << 16);
        const TIMER_PERIODIC_MODE = 1 << 17;
    }

    pub struct TDCRFlags: u32 {
        const TIMEDIVIDEBY1 = 0b1011;
        const TIMEDIVIDEBY2 = 0b0000;
        const TIMEDIVIDEBY4 = 0b0001;
        const TIMEDIVIDEBY8 = 0b0010;
        const TIMEDIVIDEBY16 = 0b0011;
        const TIMEDIVIDEBY32 = 0b1000;
        const TIMEDIVIDEBY64 = 0b1001;
        const TIMEDIVIDEBY128 = 0b1010;
    }
}
    \end{minted}

    使用 \texttt{bitflags} 宏可以清晰地定义寄存器中各个位的含义，避免直接使用魔法数字。例如，\texttt{LAPICFlags::ENAPIC} 明确表示这是启用 APIC 的标志位，而不是一个神秘的 \texttt{1 << 8}。

    同时，使用 \texttt{bit\_field} crate 的 \texttt{get\_bit} 方法检查位状态：

    \begin{minted}{rust}
fn set_icr(&mut self, value: u64) {
    unsafe {
        while self.read(0x0300).get_bit(12) {}
        self.write(0x0310, (value >> 32) as u32);
        self.write(0x0300, value as u32);
        while self.read(0x0300).get_bit(12) {}
    }
}
    \end{minted}

    ICR（Interrupt Command Register）由两个 32 位寄存器组成（0x300 和 0x310）。写入 ICR 时需要先写入高 32 位（0x310），再写入低 32 位（0x300）。第 12 位（DS，Delivery Status）表示中断是否正在发送，需要等待该位清零才能进行下一次写入。

    \item \textbf{中断系统初始化}

    在 \texttt{src/interrupt/mod.rs} 中，\texttt{init} 函数统一初始化中断系统：

    \begin{minted}{rust}
pub fn init() {
    IDT.load();

    if XApic::support() {
        let mut lapic = unsafe {
            XApic::new(physical_to_virtual(LAPIC_ADDR))
        };
        lapic.cpu_init();
        info!("APIC Initialized.");

        let mut ioapic = unsafe {
            IoApic::new(physical_to_virtual(IOAPIC_ADDR))
        };
        ioapic.disable_all();
        enable_irq(Irq::Serial0 as u8, 0);
        info!("IOAPIC Enabled.");
    } else {
        info!("APIC is not supported.");
    }
    
    info!("Interrupts Initialized.");
}
    \end{minted}

    这里需要注意，APIC 的地址（\texttt{LAPIC\_ADDR = 0xFEE00000}）是物理地址，需要通过 \texttt{physical\_to\_virtual} 函数转换为虚拟地址才能访问。这是因为内核运行在高地址空间，直接访问物理地址会导致 Page Fault。

    IOAPIC 的初始化包括禁用所有 IRQ，然后只启用需要的 IRQ（这里是串口 IRQ4）。这样可以避免收到不需要的中断。

    \item \textbf{思考题：APIC 初始化顺序}

    \textbf{【可选思考题 1】} APIC 初始化步骤必须严格按顺序执行：

    \begin{enumerate}
        \item \textbf{必须先启用 APIC（SPIV 寄存器）}，然后才能配置其他寄存器。如果先配置 LVT 或定时器寄存器，写入操作可能被忽略，因为 APIC 尚未启用。
        \item \textbf{必须先清除错误状态寄存器}，然后才能启用中断。否则残留的错误状态可能导致意外中断或中断被屏蔽。
        \item \textbf{必须先禁用 LINT0/LINT1}，然后才能发送 INIT IPI。否则可能收到意外的 NMI（Non-Maskable Interrupt）或 SMI（System Management Interrupt）。
        \item \textbf{必须先发送 INIT De-assert IPI}，然后才能启用定时器中断。否则 APIC 可能处于不一致状态，导致定时器中断不触发。
    \end{enumerate}

    如果故意交换 SPIV 启用和 LVT 配置两步，最可能出现的首个可观测故障是：LVT 配置被忽略，定时器中断不会触发，系统不会输出 "Tick!" 日志。这是因为在 APIC 未启用时，对 LVT 寄存器的写入可能不会生效。

\end{enumerate}

\section{时钟中断}

\begin{enumerate}[label={\chinese*、},labelsep=0pt, start=6]

    \item \textbf{时钟中断处理程序}

    在 \texttt{src/interrupt/clock.rs} 中，为 Timer IRQ 设置中断处理程序：

    \begin{minted}{rust}
use super::consts::*;
use x86_64::structures::idt::{InterruptDescriptorTable, InterruptStackFrame};
use core::sync::atomic::{AtomicU64, Ordering};

pub unsafe fn register_idt(idt: &mut InterruptDescriptorTable) {
    idt[Interrupts::IrqBase as u8 + Irq::Timer as u8]
        .set_handler_fn(clock_handler);
}

pub extern "x86-interrupt" fn clock_handler(_sf: InterruptStackFrame) {
    x86_64::instructions::interrupts::without_interrupts(|| {
        inc_counter();
        super::ack();
    });
}

static COUNTER: AtomicU64 = AtomicU64::new(0);

#[inline]
pub fn read_counter() -> u64 {
    COUNTER.load(Ordering::Relaxed)
}

#[inline]
pub fn inc_counter() -> u64 {
    COUNTER.fetch_add(1, Ordering::Relaxed)
}
    \end{minted}

    时钟中断处理程序的核心工作是增加计数器并发送 EOI（End Of Interrupt）信号。计数器使用 \texttt{AtomicU64} 类型，确保在并发访问时的原子性。

    \texttt{without\_interrupts} 函数在执行闭包期间临时禁用 CPU 中断。这是必要的，因为：
    \begin{itemize}
        \item 防止时钟中断嵌套：如果在中断处理过程中又收到时钟中断，会导致递归调用
        \item 确保 \texttt{ack()} 操作的完整性：EOI 信号需要在中断禁用状态下发送，以避免竞态条件
    \end{itemize}

    \item \textbf{思考题：without\_interrupts 的作用}

    \textbf{【思考题 1】} 为什么需要在 \texttt{clock\_handler} 中使用 \texttt{without\_interrupts} 函数？

    \texttt{without\_interrupts} 函数在执行闭包期间临时禁用 CPU 中断。如果不使用它，可能会发生以下情况：

    \begin{enumerate}
        \item 时钟中断处理程序执行过程中，另一个时钟中断可能嵌套触发
        \item 嵌套中断可能导致 \texttt{COUNTER} 的读写操作被中断，虽然 \texttt{AtomicU64} 本身是原子的，但 \texttt{ack()} 操作可能被中断
        \item 如果 \texttt{ack()} 被延迟，APIC 可能认为中断未被确认，导致中断丢失或重复
        \item 在更复杂的场景中（如进程调度），嵌套中断可能导致栈溢出或状态不一致
    \end{enumerate}

    需要注意的是，\texttt{without\_interrupts} 只禁用当前 CPU 的中断，不会影响其他 CPU（在多核系统中）。在单核系统中，这足以保证临界区的互斥访问。

    \item \textbf{思考题：时钟中断频率设置}

    \textbf{【思考题 2】} 时钟中断的频率设置考虑：

    \begin{itemize}
        \item \textbf{频率太快}：中断开销增大，CPU 大量时间用于处理中断，降低系统吞吐量；但进程调度更及时，响应性更好
        \item \textbf{频率太慢}：中断开销减小，但进程调度延迟增大，交互式应用响应变慢
    \end{itemize}

    在代码中，时钟频率由 TICR（初始计数值）和 TDCR（分频系数）决定：
    \begin{minted}{rust}
self.write(0x3E0, TDCRFlags::TIMEDIVIDEBY1.bits()); // 分频系数为 1
self.write(0x380, 0x20000); // 初始计数值为 0x20000
    \end{minted}

    如果想要中断频率减半，可以将初始计数值加倍：
    \begin{minted}{rust}
self.write(0x380, 0x40000); // 初始计数值加倍，频率减半
    \end{minted}

    这是因为定时器每次计数到 0 时触发中断，然后重新加载初始值。初始值越大，计数到 0 所需的时间越长，中断频率越低。

    \begin{figure}[h]
        \centering
        \includegraphics[width=0.9\textwidth]{clock_interrupt.png}
        \caption{时钟中断输出 - Tick! 日志}
    \end{figure}

    \item \textbf{加分项：线程安全的 COUNTER}

    \textbf{【加分项 2】} COUNTER 使用 \texttt{AtomicU64} 定义，能够做到线程安全：

    \begin{minted}{rust}
static COUNTER: AtomicU64 = AtomicU64::new(0);

#[inline]
pub fn read_counter() -> u64 {
    COUNTER.load(Ordering::Relaxed)
}

#[inline]
pub fn inc_counter() -> u64 {
    COUNTER.fetch_add(1, Ordering::Relaxed)
}
    \end{minted}

    \texttt{AtomicU64} 提供了原子操作：
    \begin{itemize}
        \item \texttt{load(Ordering::Relaxed)}：原子读取
        \item \texttt{fetch\_add(1, Ordering::Relaxed)}：原子增加并返回旧值
    \end{itemize}

    由于中断处理程序可能在任何时刻被触发，使用原子类型确保即使在多核环境下（虽然本实验是单核），计数器的读写也是线程安全的。\texttt{Ordering::Relaxed} 是最宽松的内存序，只保证原子性，不保证操作顺序。对于计数器这种不依赖其他内存操作的场景，\texttt{Relaxed} 序是足够的。

\end{enumerate}

\section{串口输入中断}

\begin{enumerate}[label={\chinese*、},labelsep=0pt, start=8]

    \item \textbf{串口驱动开启中断}

    在 \texttt{src/drivers/uart16550.rs} 的 \texttt{init} 函数末尾为串口设备开启中断：

    \begin{minted}{rust}
// Enable Interrupts
unsafe {
    ier_port.write(IERFlags::RECEIVE_DATA.bits());
}
    \end{minted}

    IER（Interrupt Enable Register，偏移量 1）用于启用或禁用串口的各种中断。\texttt{IERFlags::RECEIVE\_DATA} 对应第 0 位，启用接收数据中断。当串口收到数据时，会触发 IRQ4 中断。

    需要注意的是，在初始化过程中，首先禁用了所有中断（写入 0x00），完成配置后才启用需要的中断。这是为了避免在配置过程中收到意外中断。

    \item \textbf{串口中断处理程序}

    在 \texttt{src/interrupt/serial.rs} 中，为 IRQ4 Serial0 设置中断处理程序：

    \begin{minted}{rust}
use crate::drivers;

use super::consts::*;
use x86_64::structures::idt::{InterruptDescriptorTable, InterruptStackFrame};

pub unsafe fn register_idt(idt: &mut InterruptDescriptorTable) {
    idt[Interrupts::IrqBase as u8 + Irq::Serial0 as u8]
        .set_handler_fn(serial_handler);
}

pub extern "x86-interrupt" fn serial_handler(_st: InterruptStackFrame) {
    receive();
    super::ack();
}

fn receive() {
    if let Some(mut serial) = drivers::serial::get_serial() {
        if let Some(ch) = serial.receive() {
            drivers::input::push_key(ch);
        }
    }
}
    \end{minted}

    串口中断处理程序遵循 Top half \& Bottom half 原则：在中断处理中只做最少的工作（读取字符并放入缓冲区），实际的数据处理在中断之外进行。这样可以减少中断处理的时间，提高系统响应性。

    \texttt{receive} 函数通过 \texttt{get\_serial} 获取串口驱动的可变引用，然后调用 \texttt{receive} 方法读取字符。这里使用了 \texttt{Option} 类型来处理可能的空值，确保代码的健壮性。

    \item \textbf{输入缓冲区驱动}

    在 \texttt{src/drivers/input.rs} 中，使用 \texttt{crossbeam\_queue::ArrayQueue} 实现无锁输入缓冲区：

    \begin{minted}{rust}
use core::char;

use alloc::string::String;
use crossbeam_queue::ArrayQueue;
type Key = u8;

lazy_static! {
    static ref INPUT_BUF: ArrayQueue<Key> = ArrayQueue::new(128);
}

#[inline]
pub fn push_key(key: Key) {
    if INPUT_BUF.push(key).is_err() {
        warn!("Input buffer is full. Dropping key '{:?}'", key);
    }
}

#[inline]
pub fn try_pop_key() -> Option<Key> {
    INPUT_BUF.pop()
}

pub fn pop_key() -> Key {
    loop {
        if let Some(key) = try_pop_key() {
            return key;
        }
    }
}
    \end{minted}

    \texttt{ArrayQueue} 是一个无锁的、固定大小的队列，适合在中断处理程序中使用。它的 \texttt{push} 和 \texttt{pop} 操作都是原子的，不需要额外的锁。

    缓冲区大小设为 128，这是一个经验值：足够大以容纳用户快速输入的数据，又不会占用过多内存。如果缓冲区满了，\texttt{push} 会返回错误，此时输出警告日志并丢弃新数据。

    \texttt{pop\_key} 函数是一个阻塞读取函数，循环等待直到缓冲区中有数据。这个函数在用户交互时使用，确保能够获取到用户输入。

    \item \textbf{get\_line 函数实现}

    \texttt{get\_line} 函数从缓冲区阻塞读取数据，支持退格和 UTF-8 编码：

    \begin{minted}{rust}
pub fn get_line() -> String {
    let mut buf = String::with_capacity(128);
    let mut utf8_buf = [0u8; 4];
    let mut utf8_len = 0;

    // State machine for ANSI escape sequences (0: Normal, 1: ESC, 2: ESC [)
    let mut escape_state = 0;

    loop {
        let c = pop_key();
        
        // Handle ANSI escape sequences
        if escape_state == 1 {
            escape_state = if c == b'[' || c == b'O' { 2 } else { 0 };
            continue;
        } else if escape_state == 2 {
            if c >= 0x40 && c <= 0x7E { escape_state = 0; }
            continue;
        } else if c == 0x1B {
            escape_state = 1;
            utf8_len = 0;
            continue;
        }

        // Ignore unhandled control characters
        if is_ignored_ctrl(c) {
            utf8_len = 0;
            continue;
        }

        if c == 0x08 || c == 0x7f {
            if let Some(ch) = buf.pop() {
                let width = if ch.is_ascii() { 1 } else { 2 };
                for _ in 0..width {
                    print!("\x08 \x08");
                }
            }
            utf8_len = 0;
        } else if c == 0x0A || c == 0x0D {
            print!("\n");
            return buf;
        } else {
            utf8_buf[utf8_len] = c;
            utf8_len += 1;

            // Attempt to parse UTF-8 sequence
            match core::str::from_utf8(&utf8_buf[..utf8_len]) {
                Ok(s) => {
                    if let Some(ch) = s.chars().next() {
                        buf.push(ch);
                        print!("{}", ch);
                    }
                    utf8_len = 0;
                }
                Err(e) => {
                    if e.error_len().is_some() || utf8_len >= 4 {
                        utf8_len = 0;
                    }
                }
            }
        }
    }
}
    \end{minted}

    \texttt{get\_line} 函数的实现有几个关键点：

    首先，使用 \texttt{String::with\_capacity(128)} 预先分配足够的内存，避免在读取过程中频繁重新分配。

    其次，处理 ANSI 转义序列。终端发送的方向键等特殊键会产生 ANSI 转义序列（以 ESC [ 开头），这些序列需要被跳过，否则会显示为乱码。状态机用于跟踪当前的转义序列解析状态。

    第三，处理退格键（0x08 和 0x7F）。退格时需要从缓冲区删除最后一个字符，并在屏幕上擦除对应的显示。对于非 ASCII 字符（如中文、emoji），需要估计显示宽度（ASCII 为 1，其他为 2），以确保正确擦除。

    第四，UTF-8 解码。UTF-8 是一种变长编码，ASCII 字符占 1 字节，中文通常占 3 字节，emoji 可能占 4 字节。函数使用一个 4 字节的缓冲区累积 UTF-8 序列，然后尝试解析。如果解析成功，将字符添加到结果字符串；如果解析失败或超过最大长度，重置缓冲区。

    \item \textbf{思考题：串口与互斥锁}

    \textbf{【思考题 3】} 在进行 \texttt{receive} 操作的时候，为什么无法进行日志输出？

    串口驱动使用互斥锁保护共享资源。当串口中断处理程序正在读取数据时，它持有了串口的锁。如果此时尝试输出日志（也使用串口），日志系统会尝试获取同一个锁，导致死锁。

    与用户态多线程不同，中断导致的"并发访问"是强制性的：
    \begin{enumerate}
        \item 主程序持有串口锁进行日志输出
        \item 中断触发，强制打断主程序
        \item 中断处理程序尝试获取串口锁读取数据
        \item 由于锁已被持有，中断处理程序等待
        \item 但主程序被中断打断，无法释放锁
        \item 形成死锁：中断等待锁，锁持有者被中断打断
    \end{enumerate}

    这就是为什么在中断处理程序中应该做尽量少的事，避免调用可能获取锁的函数。在本实验中，\texttt{receive} 函数通过 \texttt{get\_serial} 获取串口引用，如果串口驱动使用互斥锁保护，这里可能会死锁。实际实现中，\texttt{get\_serial} 返回的是 \texttt{Option}，如果无法获取锁则返回 \texttt{None}，避免了死锁。

    \item \textbf{思考题：输入缓冲区满的情况}

    \textbf{【思考题 4】} 输入缓冲区在以下情况下会满：
    \begin{itemize}
        \item 用户快速输入大量字符，超过缓冲区容量（128 字节）
        \item 主程序长时间不读取缓冲区，中断持续接收字符
    \end{itemize}

    如果缓冲区满了，新输入的字符会被丢弃，并输出警告日志：
    \begin{minted}{rust}
if INPUT_BUF.push(key).is_err() {
    warn!("Input buffer is full. Dropping key '{:?}'", key);
}
    \end{minted}

    在实际使用中，128 字节的缓冲区通常足够大，因为用户输入速度有限。但如果主程序长时间阻塞，缓冲区可能会满。增加缓冲区大小可以缓解这个问题，但会占用更多内存。

    \item \textbf{加分项：支持中文和 emoji 输入}

    \textbf{【加分项 4】} 串口输入驱动通过 UTF-8 解码支持中文和 emoji 输入：

    在 \texttt{get\_line} 函数中，使用 UTF-8 缓冲区累积多字节字符：
    \begin{minted}{rust}
let mut utf8_buf = [0u8; 4];
let mut utf8_len = 0;

// ... 在循环中累积 UTF-8 字节
utf8_buf[utf8_len] = c;
utf8_len += 1;

// 尝试解析 UTF-8 序列
match core::str::from_utf8(&utf8_buf[..utf8_len]) {
    Ok(s) => {
        if let Some(ch) = s.chars().next() {
            buf.push(ch);
            print!("{}", ch);
        }
        utf8_len = 0;
    }
    Err(e) => {
        if e.error_len().is_some() || utf8_len >= 4 {
            utf8_len = 0;
        }
    }
}
    \end{minted}

    UTF-8 编码规则：
    \begin{itemize}
        \item 0xxxxxxx：单字节 ASCII
        \item 110xxxxx 10xxxxxx：两字节字符
        \item 1110xxxx 10xxxxxx 10xxxxxx：三字节字符（中文通常在此范围）
        \item 11110xxx 10xxxxxx 10xxxxxx 10xxxxxx：四字节字符（emoji 通常在此范围）
    \end{itemize}

    函数逐个接收字节，累积到 \texttt{utf8\_buf} 中，然后尝试解析为 UTF-8 字符串。如果解析成功，说明收到了完整的字符；如果解析失败，继续接收下一个字节。如果累积了 4 个字节仍然无法解析，说明收到了无效的 UTF-8 序列，重置缓冲区。

    退格处理也考虑了字符显示宽度（ASCII 为 1，CJK/Emoji 为 2）：
    \begin{minted}{rust}
if let Some(ch) = buf.pop() {
    let width = if ch.is_ascii() { 1 } else { 2 };
    for _ in 0..width {
        print!("\x08 \x08");
    }
}
    \end{minted}

    退格序列 \texttt{\textbackslash x08 \textbackslash x20 \textbackslash x08} 的含义是：光标左移、输出空格擦除字符、光标再左移。对于宽字符，需要重复两次这个序列。

    \begin{figure}[h]
        \centering
        \includegraphics[width=0.9\textwidth]{serial_echo_supporting_UTF-8.png}
        \caption{串口输入支持 UTF-8（中文和 emoji）}
    \end{figure}

\end{enumerate}

\section{用户交互}

\begin{enumerate}[label={\chinese*、},labelsep=0pt, start=11]

    \item \textbf{命令行交互}

    在 \texttt{src/main.rs} 中，使用 \texttt{get\_line} 函数实现用户交互：

    \begin{minted}{rust}
#![no_std]
#![no_main]

use ysos_kernel::{self as ysos, input, interrupt, print, println};

boot::entry_point!(kernel_main);

pub fn kernel_main(boot_info: &'static boot::BootInfo) -> ! {
    ysos::init(boot_info);

    loop {
        print!("> ");
        let input = input::get_line();

        match input.trim() {
            "exit" => break,
            _ => {
                println!("You said {}", input);
                println!("The counter value is {}", interrupt::clock::read_counter());
            }
        }
    };
    ysos::shutdown();
}
    \end{minted}

    主循环首先输出提示符 \texttt{"> "}，然后调用 \texttt{get\_line} 阻塞等待用户输入。用户输入一行后，检查内容：如果是 "exit"，退出循环并关闭系统；否则，回显输入并显示当前时钟计数器值。

    显示计数器值是为了验证时钟中断正常工作。如果时钟中断没有触发，计数器值将保持为 0。

    \begin{figure}[h]
        \centering
        \includegraphics[width=0.9\textwidth]{exit_command.png}
        \caption{用户交互 - exit 命令}
    \end{figure}

    \item \textbf{思考题：Triple Fault 和 Double Fault 触发}

    \textbf{【思考题 5 - Triple Fault 和 Double Fault 部分】}

    \textbf{触发 Double Fault 的方式：} 在 IDT 中未注册的处理程序触发异常，或者在异常处理程序中再次触发异常。例如，在普通异常处理程序中访问非法地址。

    \textbf{触发 Triple Fault 的方式：} Double Fault 处理程序中再次发生异常，或者 Double Fault 的 IST 栈溢出。

    开启 \texttt{intdbg} 选项可以在 QEMU 中查看调试信息：
    \begin{minted}[linenos=false]{bash}
make intdbg
    \end{minted}

    \begin{figure}[h]
        \centering
        \includegraphics[width=0.9\textwidth]{debug_using_gdb_and_tmux.png}
        \caption{使用 GDB 和 tmux 调试}
    \end{figure}

    在调试 Triple Fault 时，QEMU 会输出类似以下的信息：
    \begin{minted}[linenos=false]{text}
check_exception old: 0xffffffff new 0x8
     0: v=08 e=0000 i=0 cpl=0 IP=0008:0000000000000000 pc=0000000000000000 SP=0010:0000000000000000 CR2=0000000000000000
check_exception old: 0x8 new 0xd
     1: v=0d e=0000 i=0 cpl=0 IP=0008:0000000000000000 pc=0000000000000000 SP=0010:0000000000000000 CR2=0000000000000000
check_exception old: 0xd new 0x8
     2: v=08 e=0000 i=0 cpl=0 IP=0008:0000000000000000 pc=0000000000000000 SP=0010:0000000000000000 CR2=0000000000000000
Triple fault! Resetting...
    \end{minted}

    这里 \texttt{v=08} 表示 Double Fault（向量号 8），\texttt{v=0d} 表示 General Protection Fault（向量号 13）。可以看到异常链：Double Fault 触发 General Protection Fault，再触发 Double Fault，最终导致 Triple Fault。

\end{enumerate}

\section{实验总结}

本次实验完成了以下内容：
\begin{enumerate}
    \item 实现了 GDT 和 TSS，为 Double Fault 和 Page Fault 准备了独立的 IST 栈。IST 栈使用静态缓冲区分配在 BSS 段，确保在内核启动早期可用。
    \item 注册了全部 CPU 异常的中断处理程序（加分项 1），包括 Faults、Traps 和 Aborts 三类异常。Page Fault 和 Double Fault 使用独立的 IST 栈。
    \item 实现了 APIC 初始化，使用 bitflags 和 bit\_field 定义寄存器（加分项 3）。APIC 初始化需要严格按照特定顺序执行，否则可能导致中断不触发或系统不稳定。
    \item 实现了时钟中断，使用 AtomicU64 保证线程安全。时钟频率由 TDCR 和 TICR 寄存器共同决定。
    \item 实现了串口输入中断和输入缓冲区驱动。使用无锁队列 \texttt{ArrayQueue} 实现输入缓冲区，支持中断处理程序和主程序并发访问。
    \item 实现了支持 UTF-8 编码的用户交互（加分项 4），可以正确处理中文和 emoji 输入。
    \item 完成了所有思考题的回答，深入理解了中断处理机制、栈切换、死锁等问题。
\end{enumerate}

通过本次实验，深入了解了中断处理机制、APIC 可编程中断控制器、IDT 中断描述符表、TSS 任务状态段以及 IST 中断栈表的工作原理。掌握了中断处理程序的编写方法，理解了 Top half \& Bottom half 原则在中断处理中的应用。为后续的进程调度和系统调用实现奠定了基础。

\clearpage
\end{document}
